\section{Conclusion}
\label{sec:conclusion}

In this paper, we presented SwarmSort, a comparison of particle swarm optimization, random search and the default settings for tuning a lightweight waste detector under an equal budget of 32 proxy trainings.

With regard to RQ1, PSO did not outperform random search in any meaningful way. In \genone, which used studio data, the three arms were statistically tied (0.676, 0.672 and 0.668 \mapfifty). In \gentwo, both tuned arms were better than the defaults on the mixed test set (random search by $+0.017$ and PSO by $+0.020$), but PSO and random search could not be distinguished on the mixed, studio and real-world test sets. Their only significant difference, $+0.003$ on the India set, occurs at scores of about 0.02 and has no practical value. With regard to RQ2, the short proxy ranked the settings like a longer training ($\rho = 0.976$, $n = 8$, \genone only). The small differences it showed between the best settings, however, disappeared after full training. With regard to RQ3, the answer does depend on the domain. The gains from tuning did not carry over to real-world photographs (0.298 to 0.310), and all arms failed on the Indian DWSD photographs (0.023 to 0.027), which is partly due to the merged boxes of DWSD. Real-world training data, on the other hand, raised real-world \mapfifty from 0.038--0.096 in \genone to 0.338--0.353 in \gentwo on the same images, with a loss of 0.074 to 0.105 on studio images.

Two conclusions can be drawn from this work. First, a claim that a metaheuristic tunes a model better should be supported by the defaults and random search at an equal budget, and by confidence intervals on a test split that is used once. Had random search been left out, the \gentwo gain of PSO would have been attributed to the swarm. Second, for waste detection in real scenes, and particularly in Indian scenes, the training data are much more important than the optimizer. As discussed in Section~\ref{sec:future}, our immediate future work is to collect Indian training data with per-item boxes and to build an ensemble of the three models.
